% 决策类：E 相等 D 不同，方案 A 与方案 B 柱状图并排
\documentclass[tikz,border=4pt]{standalone}
\usepackage{ctex}
\usepackage{amsmath}
\usepackage{pgfplots}
\pgfplotsset{compat=1.18}
\usetikzlibrary{calc}
\begin{document}
\begin{tikzpicture}
\begin{axis}[
    name=planA,
    width=7.5cm, height=6cm,
    ybar,
    bar width=0.7cm,
    xmin=0.5, xmax=9.5,
    ymin=0, ymax=0.75,
    xtick={2,4,6},
    ytick={0,0.1,0.2,0.3,0.4,0.5,0.6},
    xlabel={$X_A$ 取值}, ylabel={概率},
    title={方案 A：$E_A=4$，$D_A=2.4$},
    nodes near coords,
    nodes near coords style={font=\footnotesize},
    every node near coord/.append style={/pgf/number format/.cd, fixed, precision=2},
    enlarge x limits=0.1,
]
\addplot[fill=blue!50, draw=blue!70!black] coordinates {
    (2, 0.3) (4, 0.4) (6, 0.3)
};
% 标 E_A
\draw[red, dashed, thick] (axis cs:4, 0) -- (axis cs:4, 0.65);
\node[red, font=\footnotesize, anchor=south] at (axis cs:4, 0.65) {$E_A=4$};
\end{axis}

\begin{axis}[
    at={(planA.east)}, anchor=west, xshift=1.0cm,
    width=7.5cm, height=6cm,
    ybar,
    bar width=0.7cm,
    xmin=-1.5, xmax=9.5,
    ymin=0, ymax=0.75,
    xtick={0,4,8},
    ytick={0,0.1,0.2,0.3,0.4,0.5,0.6},
    xlabel={$X_B$ 取值}, ylabel={概率},
    title={方案 B：$E_B=4$，$D_B=6.4$},
    nodes near coords,
    nodes near coords style={font=\footnotesize},
    every node near coord/.append style={/pgf/number format/.cd, fixed, precision=2},
    enlarge x limits=0.1,
]
\addplot[fill=orange!60, draw=orange!70!black] coordinates {
    (0, 0.2) (4, 0.6) (8, 0.2)
};
% 标 E_B
\draw[red, dashed, thick] (axis cs:4, 0) -- (axis cs:4, 0.65);
\node[red, font=\footnotesize, anchor=south] at (axis cs:4, 0.65) {$E_B=4$};
\end{axis}

% 底部说明（用 planA 的 below 锚定，避免复杂坐标计算）
\node[font=\small, align=center, anchor=north, yshift=-0.8cm]
  at ($(planA.south east)+(1.5cm,0)$)
  {期望相等：$E_A=E_B=4$\quad 但 方差不同：$D_A=2.4 < D_B=6.4$\\
   \footnotesize 方差越小，取值越集中，方案越稳定。};
\end{tikzpicture}
\end{document}
